\scanpage{327}\mspage{243}
% l. 1–2
4.10. \uline{Le théorème de représentation indicielle} : \uline{démonstration}.

% l. 3–5
On va d'abord reformuler le th. 4.9.3. sous forme plus générale, en termes de catégories \uline{cofibrées}. Soit
% l. 6
\[
 \mathscr D\,\deleted{\unsure{D_j}}\longrightarrow J \tag*{(4.10.1)}
\]
% l. 7–11
un foncteur cofibrant, où $J$ est équivalente à une petite catégorie, et de $\overline{\mathrm m}$ pour les catégories fibres \deleted{$\unsure{C_j|D_j}$} $\mathscr D_j$ $(j\in\operatorname{Ob}J)$. On \deleted{suppose qu'on a} \inserted{se donne} un cardinal \inserted{infini} $\unsure{\pi|\kappa}$ \deleted{tel que}
\[
 \deleted{\operatorname{Card}\unsure{\mathrm{réd}|\mathrm{Fl}\,\mathrm{de}|\mathrm{Flct}}\,J\leq\pi}. \tag*{\deleted{(4.10.2)}}
\]
% l. 12–13
Soit
\[
 \unsure{M|\mathscr M}_j\deleted{\unsure{j}}=\operatorname{Ind}_{\pi}(\mathscr D_j)
 \qquad(j\in\operatorname{Ob}J), \tag*{(4.10.2)}
\]
% l. 14–18
alors les $M_j$ sont les fibres d'une catégorie cofibrée \inserted{$\unsure{\mathscr N|\mathscr M}$} sur $J$, que je désigne aussi par $\operatorname{Ind}_{\pi}(\mathscr D/J)$, donc on pose
\[
 \mathscr N=\operatorname{Ind}_{\pi}(\mathscr D/J). \tag*{(\unsure{4.10.3})}
\]
% l. 19–23
On \deleted{\unsure{peut}} a donc une inclusion \inserted{pl. fid.} de catégories cofibrées sur $J$
\[
 \mathscr D\hookrightarrow\mathscr N,
\]
d'où une \unsure{inclusion|induction}
\[
 \underline{\Gamma}(\mathscr D/J)\longrightarrow\underline{\Gamma}(\mathscr N/J). \tag*{(4.10.4)}
\]
% archive sheet number: 326
